\section{Ортогональные операторы (изометрии)}

\begin{definition}
Оператор $\mathcal Q$ евклидова пространства называется \emph{ортогональным}, если он сохраняет
скалярное произведение: $\scal{\mathcal Qx}{\mathcal Qy}=\scal xy$ для всех $x,y$.
\end{definition}

\begin{theorem}[критерии ортогональности]
Следующие условия равносильны:
\begin{enumerate}[label=\arabic*),leftmargin=*,nosep]
  \item $\mathcal Q$ сохраняет скалярное произведение;
  \item $\mathcal Q$ сохраняет норму: $\norm{\mathcal Qx}=\norm x$ для всех $x$;
  \item $\mathcal Q^*\mathcal Q=\mathcal I$, т.е.\ $\mathcal Q^*=\mathcal Q^{-1}$ (в ОНБ
        $Q^{T}Q=I$).
\end{enumerate}
\end{theorem}

\begin{proof}
$1\Rightarrow2$: при $y=x$ получаем $\norm{\mathcal Qx}^2=\norm x^2$. $2\Rightarrow1$: по формуле
поляризации $\scal xy=\tfrac14(\norm{x+y}^2-\norm{x-y}^2)$ сохранение нормы влечёт сохранение
произведения. $1\Leftrightarrow3$: $\scal{\mathcal Qx}{\mathcal Qy}=\scal{x}{\mathcal Q^*\mathcal Qy}$,
и это равно $\scal xy$ для всех $x,y$ тогда и только тогда, когда $\mathcal Q^*\mathcal Q=\mathcal I$.
\end{proof}

\section{Свойства ортогональных матриц}

\begin{property}\leavevmode
\begin{enumerate}[label=\arabic*),leftmargin=*,nosep]
  \item Столбцы (и строки) ортогональной матрицы образуют ОНБ ($Q^{T}Q=I$ — это и есть условие
        $\scal{q_i}{q_j}=\delta_{ij}$).
  \item $\det\mathcal Q=\pm1$.
  \item Все собственные значения по модулю равны $1$; вещественные собственные значения — только
        $\pm1$.
\end{enumerate}
\end{property}

\begin{proof}
2) Из $Q^{T}Q=I$: $(\det Q)^2=\det Q^{T}\det Q=\det I=1$, значит $\det Q=\pm1$.\;
3) Если $\mathcal Qx=\lambda x$ ($x\ne0$), то $\norm x=\norm{\mathcal Qx}=\abs\lambda\,\norm x$,
откуда $\abs\lambda=1$; для вещественного $\lambda$ это $\pm1$.
\end{proof}

\section{Классификация ортогональных операторов}

\begin{definition}
Ортогональный оператор \emph{собственный} (вращение), если $\det\mathcal Q=+1$, и
\emph{несобственный} (с отражением), если $\det\mathcal Q=-1$.
\end{definition}

\begin{remark}[геометрия в $\R^2$ и $\R^3$]
\begin{itemize}[leftmargin=*,nosep]
  \item $\R^2$: при $\det=+1$ матрица есть поворот
        $\begin{psmallmatrix}\cos\varphi&-\sin\varphi\\ \sin\varphi&\cos\varphi\end{psmallmatrix}$;
        при $\det=-1$ — отражение относительно прямой.
  \item $\R^3$: при $\det=+1$ — поворот вокруг некоторой оси (она отвечает собственному значению
        $+1$); при $\det=-1$ — поворот, скомбинированный с отражением (зеркальный поворот).
\end{itemize}
В подходящем ОНБ любой ортогональный оператор блочно-диагонален: блоки $2\times2$ — повороты на углы
$\varphi_k$, плюс отдельные значения $\pm1$ на диагонали.
\end{remark}
